\usepackage[T1]{fontenc}
\usepackage[utf8]{inputenc}
\usepackage{lmodern}
\usepackage{microtype}
\usepackage[a4paper,margin=1in]{geometry}

\usepackage{amsmath,amssymb,amsthm,mathtools}
\usepackage{bm}
\usepackage{booktabs}
\usepackage{array}
\usepackage{tabularx}
\usepackage{longtable}
\usepackage{enumitem}
\usepackage{xcolor}
\usepackage{graphicx}
\usepackage{url}
\usepackage{natbib}
\usepackage{hyperref}
\usepackage[nameinlink,capitalise,noabbrev]{cleveref}

\definecolor{deepblue}{HTML}{173B57}
\definecolor{softblue}{HTML}{EAF3F8}
\definecolor{deepgreen}{HTML}{235D3A}
\definecolor{softgreen}{HTML}{EAF5EE}
\definecolor{deeporange}{HTML}{9C4E00}
\definecolor{softorange}{HTML}{FFF1DF}
\definecolor{draftred}{HTML}{8F1D1D}
\definecolor{draftbg}{HTML}{FFF4F4}
\definecolor{openpurple}{HTML}{5B3F8C}
\definecolor{softpurple}{HTML}{F2ECFA}
\definecolor{graytext}{HTML}{4B5563}

\hypersetup{
  colorlinks=true,
  linkcolor=deepblue,
  citecolor=deepgreen,
  urlcolor=deepblue,
  pdftitle={Nuclear-Norm Best Approximation at Rank-Deficient Matrix-Line Crossings},
  pdfauthor={Anonymous manuscript}
}

\setlength{\parindent}{0pt}
\setlength{\parskip}{0.55em}
\setlist[itemize]{topsep=0.35em,itemsep=0.2em}
\setlist[enumerate]{topsep=0.35em,itemsep=0.25em}
\renewcommand{\arraystretch}{1.15}

\newtheoremstyle{reader}
  {0.8em}{0.8em}
  {\itshape}{}
  {\bfseries\color{deepblue}}{.}
  {0.5em}{}
\theoremstyle{reader}
\newtheorem{theorem}{Theorem}[section]
\newtheorem{proposition}[theorem]{Proposition}
\newtheorem{lemma}[theorem]{Lemma}
\newtheorem{corollary}[theorem]{Corollary}

\theoremstyle{definition}
\newtheorem{definition}[theorem]{Definition}
\newtheorem{assumption}[theorem]{Assumption}
\newtheorem{example}[theorem]{Example}

\theoremstyle{remark}
\newtheorem{remark}[theorem]{Remark}

\newenvironment{proofidea}
  {\par\smallskip\begin{center}\begin{minipage}{0.94\textwidth}
   {\color{deepblue}\hrule}\vspace{0.35em}
   \textbf{\color{deepblue}Proof idea.}\ \color{black}\ignorespaces}
  {\par\vspace{0.35em}{\color{deepblue}\hrule}
   \end{minipage}\end{center}\smallskip}

\newenvironment{intuition}
  {\par\smallskip\begin{center}\begin{minipage}{0.94\textwidth}
   {\color{deepgreen}\hrule}\vspace{0.35em}
   \textbf{\color{deepgreen}What this result does.}\ \color{black}\ignorespaces}
  {\par\vspace{0.35em}{\color{deepgreen}\hrule}
   \end{minipage}\end{center}\smallskip}

\newenvironment{boundary}
  {\par\smallskip\begin{center}\begin{minipage}{0.94\textwidth}
   {\color{deeporange}\hrule}\vspace{0.35em}
   \textbf{\color{deeporange}Boundary of the claim.}\ \color{black}\ignorespaces}
  {\par\vspace{0.35em}{\color{deeporange}\hrule}
   \end{minipage}\end{center}\smallskip}

\newenvironment{openproblem}
  {\par\smallskip\begin{center}\begin{minipage}{0.94\textwidth}
   {\color{openpurple}\hrule}\vspace{0.35em}
   \textbf{\color{openpurple}Open problem.}\ \color{black}\ignorespaces}
  {\par\vspace{0.35em}{\color{openpurple}\hrule}
   \end{minipage}\end{center}\smallskip}

\newcommand{\proved}{\textcolor{deepgreen}{\textbf{PROVED}}}
\newcommand{\localproved}{\textcolor{deeporange}{\textbf{PROVED UNDER LOCAL ASSUMPTIONS}}}
\newcommand{\observed}{\textcolor{deepblue}{\textbf{NUMERICALLY OBSERVED}}}
\newcommand{\openstatus}{\textcolor{openpurple}{\textbf{OPEN}}}

\newcommand{\ip}[2]{\left\langle #1,#2\right\rangle}
\newcommand{\R}{\mathbb{R}}
\newcommand{\norm}[1]{\left\lVert #1\right\rVert}
\newcommand{\snorm}[1]{\left\lVert #1\right\rVert_{2}}
\newcommand{\fnorm}[1]{\left\lVert #1\right\rVert_{\mathrm F}}
\newcommand{\nnorm}[1]{\left\lVert #1\right\rVert_{*}}
\newcommand{\argmin}{\operatorname*{arg\,min}}
\newcommand{\argmax}{\operatorname*{arg\,max}}
\newcommand{\diag}{\operatorname{diag}}
\newcommand{\rank}{\operatorname{rank}}
\newcommand{\range}{\operatorname{range}}
\newcommand{\Retr}{\operatorname{Retr}}
\newcommand{\Gap}{\operatorname{Gap}}
\newcommand{\Sc}{S_{\mathrm c}}
\newcommand{\Kcan}{K_{\mathrm{can}}}
\newcommand{\Pstar}{\mathcal P^\star}
\newcommand{\Ctheta}{\mathcal C_{\Theta}}

\DeclareMathOperator{\tr}{tr}
\DeclareMathOperator{\dist}{dist}
\DeclareMathOperator{\sign}{sign}

\crefname{theorem}{theorem}{theorems}
\crefname{proposition}{proposition}{propositions}
\crefname{lemma}{lemma}{lemmas}
\crefname{corollary}{corollary}{corollaries}
\crefname{definition}{definition}{definitions}
\crefname{assumption}{assumption}{assumptions}
\crefname{example}{example}{examples}
\crefname{remark}{remark}{remarks}
